\subsection[Basis for P\_mV\^\{otimes m\}]{Basis for $\boldsymbol{P_mV^{\otimes m}}$}

Before proving Theorem \ref{Main}, we will write a basis for the symmetric projection $P_mV^{\otimes m}$.

Let $ I_0,\dots,I_N $ be the canonical basis of $ V $, and define the action on $ V $ by 	
\begin{gather}
\hat{e}_{+,i}I_j=\delta_{ij}I_{j-1},\qquad
\hat{e}_{-,i}I_{j-1}=\delta_{ij}I_j,\qquad
q^{\pm \epsilon_i/2}I_j=q^{\pm \delta_ij/2}I_j.\label{21}
\end{gather}
For $\mu \in \mathcal{B}_m^{(N)}$, define the vector
\begin{equation*}%\label{key}
v\zo =I_0^{\otimes \mu_0} \otimes \dots \otimes I_N^{\otimes \mu_N}
\end{equation*}
Define
\begin{equation*}
M(\zz) = \sum_{\sigma \in D^{\mu}}q^{-\inv (\sigma)} \sigma (v(\mu)).
\end{equation*}
Here, as before, $S_m $ acts on $ V^{\otimes m} $ by permuting the order.
By an abuse of notation, for $\ii \in \mathcal{W}_m$ we define
\[
v(\ii) = v\big(\mu^{(N)}(\ii)\big), \qquad M(\ii) = M\big(\mu^{(N)}(\ii)\big),
\]
where $N \geq i_m$. Note that these definitions do not depend on the value of $N$.
We briefly note the identity
\begin{equation}\label{Brief}
e^{\sigma(a)}_{\pm,i} ( \sigma(v(\mu)) ) = \sigma\big( e^{(a)}_{\pm,i} ( v(\mu) ) \big) \qquad \text{for all } \sigma \in S_m.
\end{equation}

We now show that the set $ \{M(\mu)\}_{\mu\in \mathcal{B}^{(N)}_m} $ gives a basis for $ P_{m}V^{\otimes m} $. A more general statement appeared in~\cite{KIMRN} with a more complicated proof, but the expression here is more convenient for calculations.
\begin{Theorem}
	The set $ \{M(\mu)\}_{\mu\in \mathcal{B}^{(N)}_m} $ gives a basis for $ P_{m}V^{\otimes m} $.
\end{Theorem}
\begin{proof}
Note that $ \big|\mathcal{B}^{(N)}_m\big|= \dimension P_mV^{\otimes m} $ and $ \{M(\mu)\}_{\mu\in \mathcal{B}^{(N)}_m} $ is a linearly independent set, so it suffices to show that $ M(\mu) \in P_mV^{\otimes m} $ for all $ \mu \in \mathcal{B}^{(N)}_m $. To show that every $ M(\mu) $ is in $ P_{m}V^{\otimes m} $, it suffices to show that \begin{equation} \label{hati}
\Delta^{(m)}(\hat{e}_{-,i}) (M(\mu))=q^{\frac{1}{2}(\mu_{i-1}+\mu_i-1)} \big(1+q^{-2}+q^{-4}+\dots+q^{-2\mu_i}\big) \cdot M\big(\mu-\hat{i}\big),
\end{equation}
where $ \hat{i}=(0,\dots,0,1,-1,0,\dots,0) $ with the $ -1 $ occurring in the $ i $th position. By~\eqref{qerel} and~\eqref{21}, $ \Delta^{(m)} (\hat{e}_{-,i})(M(\mu)) $ is in the span of $ \{\tau(v(\mu-\hat{i}))\}_{\tau\in S_m}$. Let $ A(\tau) $ be the coefficients in the expansion
\[
 \Delta^{(m)} (\hat{e}_{-,i})(M(\mu)) = \sum_{\tau \in D^{\mu-\hat{i}}} A(\tau)\tau \big(v\big(\mu-\hat{i}\big)\big).
\]
It suffices to show that $A(\tau)= q^{\frac{1}{2}(\mu_{i-1}+\mu_i-1)}\big(1+q^{-2}+q^{-4}+\dots+q^{-2\mu_i}\big) q^{-\inv(\tau)}$ for all $\tau \in D^{\mu-\hat{i}}$. In fact, we show something stronger: for all $\tau \in D^{\mu-\hat{i}}$, there exist elements $\sigma^{(0)},\dots, \sigma^{(\mu_i)}\in D^{\mu}$ such that for $0 \leq j \leq \mu_i$,
\[
\hat{e}_{-,i}^{(a_j)} \big( q^{-\inv(\sigma^{(j)})} \sigma^{(j)}(v(\mu)) \big) = q^{\frac{1}{2}(\mu_{i-1}+\mu_i-1)} q^{-2j} q^{-\inv(\tau)} \tau\big(v\big(\mu- \hat{i}\big)\big).
\]

We will proceed by induction on the value of $\inv(\tau)$, using Lemma \ref{Gene}. The base case is when~$\tau$ is the identity permutation. Then it is straightforward to check that
\[
\sigma^{(j)} = (\mu_{[0,i-1]} \ \ \mu_{[0,i-1]}+1 \ \ \dots \ \ \mu_{[0,i-1]}+j),
\]
where $\mu_{[0,i-1]}= \mu_0 + \dots + \mu_{i-1}$, satisfies the necessary conditions.

Now fix $\tau \in D^{\mu - \hat{i}}$, and suppose that the induction hypothesis holds for $\tau$. Suppose that $\tilde{\tau} \in D^{\mu - \hat{i}}$ satisfies $\tilde{\tau}=s\tau$ for some transposition $s$, and assume without loss of generality that $\inv(\tilde{\tau}) = \inv(\tau) + 1$. Define
\[
\tilde{\sigma}^{(j)}
=
\begin{cases}
s \sigma^{(j)}, & \text{if } s \sigma^{(j)} \in D^{\mu}, \\
\sigma^{(j)}, & \text{else. }
\end{cases}
\]
We now aim to prove that
\begin{equation}\label{Aim}
\hat{e}^{(s(a_j))}_{-,i}\big(q^{-\inv(\tilde{\sigma}^{(j)})} \tilde{\sigma}^{(j)}(v(\mu)) \big) = q^{\frac{1}{2}(\mu_{i-1}+\mu_i-1)} q^{-2j} q^{-\inv(\tilde{\tau})} \tilde{\tau}\big(v\big(\mu- \hat{i}\big)\big).
\end{equation}
If $\tilde{\sigma}^{(j)} = s\sigma^{(j)}$, then \eqref{Aim} follows from \eqref{Brief} and the induction hypothesis. Now assume that $\tilde{\sigma}^{(j)} = \sigma^{(j)}$. Then $s= (a_j-1 \ \ a_j)$ and
\[
\tilde{\sigma}^{(j)} ( v(\mu) ) = \underbrace{I_* \otimes \dots \otimes I_* }_{a_j-2} \otimes I_{i-1} \otimes I_{i-1} \otimes \underbrace{I_* \otimes \dots \otimes I_*}_{m-a_j}.
\]
Using that
\[
 P \big( \big(\hat{e}_{i,-} \otimes q^{-h_i/2}\big)(I_{i-1} \otimes I_{i-1}) \big) = q^{-1}\big(q^{h_i/2} \otimes \hat{e}_{i,-}\big)(I_{i-1} \otimes I_{i-1}) ,
\]
where $P(x\otimes y)=y\otimes x$ is the permutation operator, we have that
\[
\hat{e}^{(s(a_j))}_{-,i}\big(q^{-\inv(\tilde{\sigma}^{(j)})} \tilde{\sigma}^{(j)}(v(\mu)) \big) = q^{-1}\hat{e}^{(a_j)}_{-,i}\big(q^{-\inv({\sigma}^{(j)})} {\sigma}^{(j)}(v(\mu)) \big).
\]
And now \eqref{Aim} follows from the induction hypothesis.
\end{proof}
\begin{Remark}\label{ExRep}
	For all $ \mu \in \mathcal{B}^{(N)}_m $, let $ \tilde{M}(\mu)=c(\mu)M(\mu) $ where $ c(\mu) $ is defined inductively by
	\begin{equation*}%\label{key}
	c(m,0,\dots,0)=1, \qquad
	\frac{c(\mu-i)}{c(\mu)}=\frac{q^{\frac{\mu_{i-1}+\mu_i+1}{2}}\big(1-q^{-2\mu_i}\big)}{q^{\mu_i+1}-q^{-\mu_i-1}} ,
	\end{equation*} then equation \eqref{hati} is equivalent to
\begin{equation*}%\label{key}
	\Delta^{(m)} (\hat{e}_{-,i}) \tilde{M}(\mu) =\frac{q^{\mu_i+1}-q^{-\mu_i-1}}{q-q^{-1}} \tilde{M}\big(\mu-\hat{i}\big).
	\end{equation*} In fact, one can show that
	\begin{gather*}
		\Delta^{(m)} (\hat{e}_{+,i}) \tilde{M}(\mu) =\frac{q^{\mu_{i-1}+1}-q^{-\mu_{i-1}-1}}{q-q^{-1}} \tilde{M}\big(\mu+\hat{i}\big),\\
		\Delta^{(m)} (\hat{e}_{-,i}) \tilde{M}(\mu) =\frac{q^{\mu_i+1}-q^{-\mu_i-1}}{q-q^{-1}} \tilde{M}\big(\mu-\hat{i}\big),\\
	 \Delta^{(m)} \big(q^{\pm h_i/2}\big) \tilde{M}(\mu) =q^{\pm\frac{1}{2}(\mu_{i-1}-\mu_i)} \tilde{M}(\mu)
	\end{gather*} defines a representation on $ P_mV^{\otimes m} $. This is equivalent to the representation in \cite[equation~(3)]{KMMO} and \cite[Lemma~3.1]{KIMRN}.
\end{Remark}
\begin{Corollary}\label{Core} Let $\ii,\jj\in\mathcal{W}_m$. For any $\tau,\sigma \in D_{\ii}$ and any $\zeta \in D_{\jj}$,
\[
\big(q^{\inv(\tau)}e_{\zeta(\jj) \tau(\ii)} - q^{\inv(\sigma)}e_{\zeta(\jj) \sigma(\ii)} \big) \big|_{P_{ m}^+V^{\otimes m} } =0.
\]
	
	Furthermore,
\[
\bigg( \sum_{\tau \in D_{\jj}} q^{-d_{\ii}(\tau)}e_{\tau(\jj)\tau(\ii)} \bigg) M(\ii) = M(\jj) .
\]
\end{Corollary}
\begin{proof}
For any $\mu \in \mathcal{B}^{(N)}_m$ not equal to $\mu(\ii)$, it is straightforward that
\[
\big(q^{\inv(\tau)}e_{\zeta(\jj) \tau(\ii)} - q^{\inv(\sigma)}e_{\zeta(\jj) \sigma(\ii)} \big) M(\mu)=0-0=0.
\]
On the other hand,
\begin{gather*}
\begin{split}
&\big(q^{\inv(\tau)}e_{\zeta(\jj) \tau(\ii)} - q^{\inv(\sigma)}e_{\zeta(\jj) \sigma(\ii)} \big) M(\mu(\ii))\\
&\qquad = q^{\inv(\tau)}e_{\zeta(\jj) \tau(\ii)} \big(q^{-\inv(\tau)} \tau(v(\mu(\ii)) )\big) - q^{\inv(\sigma)}e_{\zeta(\jj) \sigma(\ii)} \big(q^{-\inv(\sigma)} \sigma(v(\sigma(\ii)) )\big) \\
&\qquad = \zeta(v(\mu(\jj))) - \zeta(v(\mu(\jj)))=0.
\end{split}
\end{gather*}

For the second statement, we use that
\[
e_{\tau(\jj)\tau(\ii)}\sigma(v(\ii))
=
\begin{cases}
\tau(v(\jj)), & \tau \sim_{\ii} \sigma, \\
0, & \text{else.}
\end{cases}
\]
in order to show
\begin{align*}
\bigg( \sum_{\tau \in D_{\jj}} q^{-d_{\ii}(\tau)}e_{\tau(\jj)\tau(\ii)} \bigg) \sum_{\sigma \in D_{\ii}} q^{-\inv(\sigma)} \sigma(v(\ii)) &= \sum_{\substack{\tau \in D_{\jj},\sigma \in D_{\ii} \\ \tau\sim_{\ii} \sigma}}q^{-d_{\ii}(\tau)+ \inv(\sigma)} e_{\tau(\jj)\tau(\ii)}\sigma(v(\ii))\\
&= \sum_{\tau \in D_{\jj}} q^{-\inv(\tau)} \tau(v(\jj)).\tag*{\qed}
\end{align*}\renewcommand{\qed}{}
\end{proof}
